\section{Related Work}
\label{sec:related}

\paragraph{Robustness to common corruptions.}
ImageNet-C~\cite{hendrycks2019benchmarking} established severity-controlled common corruptions (blur, noise, weather, digital artifacts) as the standard probe of real-world robustness, and a long line of work measures and improves \emph{task accuracy} under these corruptions. For VLMs, corruption robustness has likewise been studied for captioning, VQA, and reasoning utility. \TODO{cite 2--3 VLM corruption-robustness benchmarks (e.g., benchmarks evaluating MLLMs under ImageNet-C-style corruptions) --- verify exact refs}. What this literature does not ask is whether \emph{safety behavior} survives the same corruptions; utility and safety are implicitly assumed to degrade together. Our central finding --- that they decouple --- shows this assumption is false.

\paragraph{Multimodal safety benchmarks and attacks.}
Multimodal jailbreaks began with attacks that smuggle harmful \emph{text} through the visual channel: FigStep renders harmful instructions as typographic images~\cite{gong2023figstep}, and MM-SafetyBench pairs queries with diffusion-generated and typographic images~\cite{liu2024mmsafetybench}. A second generation of benchmarks moved the harm into \emph{image semantics}: SIUO curates cases where safe text plus a safe-looking image jointly imply an unsafe request~\cite{siuo2024}, VLSBench isolates risk that is visible only in the image~\cite{hu2024vlsbench}, SPA-VL provides preference data over harmful multimodal scenes~\cite{zhang2024spavl}, HoliSafe taxonomizes safe/unsafe image--text combinations~\cite{holisafe2025}, and MSSBench probes situational safety~\cite{zhou2024mssbench}. Adversarial-image jailbreaks~\cite{qi2024visual} optimize pixel perturbations to elicit harm. Our work is orthogonal to all of these: we hold the benchmarks fixed and ask what \emph{non-adversarial, ubiquitous} degradations do to their outcomes --- and we show the answer depends systematically on which of these benchmark families the harm signal lives in.

\paragraph{Safety alignment and defenses for VLMs.}
Safety fine-tuning approaches such as VLGuard~\cite{zong2024vlguard} and reasoning-based safety tuning \TODO{cite TIS / MSR-Align --- verify exact refs and venues} substantially reduce attack success on clean benchmarks, while system-prompt and inference-time defenses offer training-free alternatives. Over-refusal benchmarks such as XSTest~\cite{rottger2024xstest} measure the tax these defenses impose on benign queries. Prior evaluations of these defenses are, again, clean-image evaluations; we show that corruption partially re-opens the very vulnerability the tuning closed, and that making either the tuning data or the system prompt corruption-aware changes robustness in measurably different ways.

\paragraph{Reasoning VLMs and safety.}
Chain-of-thought VLMs --- LLaVA-CoT~\cite{xu2024llavacot}, LlamaV-o1~\cite{thawakar2025llamavo1}, R1-Onevision~\cite{yang2025r1onevision} --- claim stronger visual reasoning; whether reasoning helps or hurts safety is contested. \TODO{cite 1--2 works on reasoning-model safety (e.g., safety of R1-style reasoning)}. By judging the reasoning trace and the final answer separately, we quantify a specific failure shape: models frequently articulate the risk in their reasoning yet still comply in their answer, and additional test-time thinking does not buy safety.
